\chapter{Estandarización de la Llave Maestra Comunal}

\section{Criterio de llave principal}

El proyecto resuelve la integración por \texttt{codigo\_comuna}. Esta decisión se sustenta en evidencia directa de las fuentes: los nombres de comuna cambian entre escritura completamente en mayúsculas, escritura capitalizada y variantes con o sin tildes. Por esa razón, \texttt{nombre\_comuna} se mantiene como apoyo descriptivo y de verificación, pero no como llave de merge.

\section{Base maestra y universo final}

La base maestra se encuentra en \texttt{data/raw/dim\_comuna\_base.csv}. El reporte de homologación confirma 32 filas, unicidad de \texttt{codigo\_comuna} y \texttt{nombre\_comuna}, provincia única \texttt{SANTIAGO} y región única \texttt{METROPOLITANA DE SANTIAGO}. Ese archivo define el universo final del proyecto y se utiliza para acotar cobertura durante transformación e integración.

\section{Resultados de homologación}

Dentro del universo final, las fuentes A y B presentan 26 coincidencias exactas y 6 coincidencias resueltas por normalización. Las fuentes C y D presentan 32 coincidencias resueltas por normalización y no dejan pendientes manuales para las comunas de interés. En todos los casos, el criterio operativo es filtrar por \texttt{codigo\_comuna} y luego sobreescribir \texttt{nombre\_comuna} con la versión estandarizada de la dimensión maestra.

\begin{table}[htbp]
\centering
\footnotesize
\caption{Resumen por fuente de la homologación dentro del universo final.}
\begin{tabular}{|c|r|r|r|r|}
\hline
Fuente & Exactas & Por normalización & Manual & Pendientes \\
\hline
A & 26 & 6 & 0 & 0 \\
B & 26 & 6 & 0 & 0 \\
C & 0 & 32 & 0 & 0 \\
D & 0 & 32 & 0 & 0 \\
\hline
\end{tabular}
\end{table}


\section{Ejemplos reales}

\begin{table}[htbp]
\centering
\footnotesize
\caption{Ejemplos reales de homologación resuelta por normalización automática.}
\begin{tabular}{|c|c|p{3.0cm}|p{3.0cm}|p{4.5cm}|}
\hline
Fuente & Código & Nombre base & Nombre observado & Resolución \\
\hline
A & 13104 & CONCHALI & CONCHALÍ & coincidencia por normalizacion \\
A & 13106 & ESTACION CENTRAL & ESTACIÓN CENTRAL & coincidencia por normalizacion \\
A & 13119 & MAIPU & MAIPÚ & coincidencia por normalizacion \\
A & 13122 & PEÑALOLEN & PEÑALOLÉN & coincidencia por normalizacion \\
A & 13129 & SAN JOAQUIN & SAN JOAQUÍN & coincidencia por normalizacion \\
A & 13131 & SAN RAMON & SAN RAMÓN & coincidencia por normalizacion \\
B & 13104 & CONCHALI & CONCHALÍ & coincidencia por normalizacion \\
B & 13106 & ESTACION CENTRAL & ESTACIÓN CENTRAL & coincidencia por normalizacion \\
\hline
\end{tabular}
\end{table}


\section{Controles de unicidad y cobertura}

La homologación no deja casos pendientes dentro de las 32 comunas del proyecto. El dataset final validado conserva una sola fila por \texttt{codigo\_comuna}, sin faltantes ni extras respecto a \texttt{dim\_comuna\_base.csv}. Esta consistencia se vuelve a confirmar en la validación final y en la carga a SQLite.
